%! AP Statistics — Unit 2, Lesson 2.5 — Guided Notes
\documentclass[10pt]{article}
\usepackage{apstats-article}
\usepackage{apstats-boxes}

\begin{document}

\pageheader{Unit 2, Lesson 2.5}{Guided Notes}
\namedateperiod

\vspace{0.1in}

\begin{objectivebox}
\small By the end of this lesson, I will be able to\dots\\[4pt]
\writeline\\[1pt]
\writeline\\[1pt]
\writeline
\end{objectivebox}

\vspace{0.08in}

\begin{vocabbox}
\termblanklong{Correlation coefficient ($r$)}
\termblanklong{Positive correlation}
\termblanklong{Negative correlation}
\termblanklong{No linear association}
\termblanklong{Lurking variable}
\end{vocabbox}

\vspace{0.08in}

\begin{hookbox}
\small\textbf{Opening Question:} Your teacher displays four scatterplots.

\vspace{6pt}
Rank them from \textit{strongest} linear association to \textit{weakest} and explain.\\[8pt]
\writeline\\[1pt]\writeline
\vspace{4pt}
\textbf{Key idea:} A single number called the \blank{5cm} lets us compare strengths precisely.
\end{hookbox}

\vspace{0.08in}

\begin{notesbox}{Step 1 --- What Is the Correlation Coefficient? (DAT-1.B)}
\small

\noindent The \textbf{correlation coefficient} $r$ measures the \blank{3cm} and \blank{3cm}
of the \blank{3cm} association between two quantitative variables.

\vspace{8pt}
\textbf{Properties of $r$:}
\begin{itemize}[leftmargin=1.5em, itemsep=3pt]
  \item $r$ is always between \blank{1cm} and \blank{1cm}, inclusive.
  \item The \blank{1.5cm} (positive/negative) of $r$ gives the \blank{2.5cm} of the association.
  \item The \blank{2cm} (size) of $r$ gives the \blank{2cm} of the linear association.
  \item $r$ is \blank{2.5cm} --- it does not change when units change.
  \item Swapping $x$ and $y$ \blank{4cm} the value of $r$.
\end{itemize}

\vspace{8pt}
\textbf{The formula (for reference):}
\[
  r = \frac{1}{n-1}\sum_{i=1}^{n}\left(\frac{x_i - \bar{x}}{s_x}\right)\!\left(\frac{y_i - \bar{y}}{s_y}\right)
\]
Each term is the product of two \blank{3cm} --- that is why $r$ is unit-free.
In practice, we use \blank{3cm}.

\vspace{8pt}
\textbf{Correlation strength guidelines (approximate):}

\vspace{4pt}
\renewcommand{\arraystretch}{1.4}
\begin{tabularx}{\linewidth}{@{} >{\bfseries}C{2.5cm} X @{}}
\rowcolor{sky}
$|r|$ range & Interpretation \\
\midrule
$0.9$ to $1.0$ & \blank{6cm} linear association \\
$0.7$ to $0.9$ & \blank{6cm} linear association \\
$0.5$ to $0.7$ & \blank{6cm} linear association \\
$0$ to $0.5$   & \blank{6cm} linear association \\
\end{tabularx}
\end{notesbox}

\vspace{0.08in}

\begin{notesbox}{Step 2 --- Interpreting $r$ in Context (DAT-1.C)}
\small

\noindent \textbf{Interpretation template:}\\[6pt]
``There is a \blank{3cm} \blank{3.5cm} \textbf{linear} association between
[explanatory variable] and [response variable] for these [individuals] ($r = \blank{1.5cm}$).''

\vspace{8pt}
\textbf{Class example --- Hours of sleep ($x$) vs.\ reaction time in ms ($y$):}

\vspace{4pt}
\renewcommand{\arraystretch}{1.4}
\begin{tabularx}{\linewidth}{@{} l *{8}{C{1cm}} @{}}
\rowcolor{sky}
\textbf{Sleep (hr)} & 5 & 6 & 7 & 7 & 8 & 8 & 9 & 9 \\
\textbf{RT (ms)}    & 310 & 285 & 260 & 250 & 235 & 220 & 205 & 195 \\
\end{tabularx}

\vspace{6pt}
Technology gives $r = \blank{1.5cm}$.

\vspace{6pt}
Write the full interpretation in context:\\[6pt]
\writeline\\[1pt]\writeline\\[1pt]\writeline

\vspace{6pt}
\textbf{Caution --- $r$ near $\pm1$ does NOT guarantee a linear model is appropriate.}\\
A dataset following a perfect \blank{3cm} pattern can have $r$ close to $1$ if the
curve rises consistently. Always check the \blank{3cm} before reporting $r$.
\end{notesbox}

\vspace{0.08in}

\begin{notesbox}{Step 3 --- Correlation Does Not Imply Causation (DAT-1.C.2)}
\small

\noindent A strong $r$ tells us two variables are \blank{3cm} --- it does \emph{not} tell
us that one \blank{3cm} the other.

\vspace{6pt}
\textbf{Classic example:} Cities with more \blank{4cm} tend to have more \blank{4cm}.
Does that mean one causes the other? No --- both are caused by a \blank{3.5cm} variable: \blank{4cm}.

\vspace{6pt}
A \textbf{lurking variable} is a variable that is \blank{5cm} and may explain the
relationship between $x$ and $y$.

\vspace{6pt}
\textbf{Practice --- For each pair, identify a possible lurking variable:}
\begin{enumerate}[leftmargin=1.5em, itemsep=4pt]
  \item Ice cream sales and drowning deaths (strong positive): \blank{6cm}
  \item Shoe size and reading level in elementary school (positive): \blank{6cm}
\end{enumerate}
\end{notesbox}

\vspace{0.08in}

\begin{spiralbox}
\small
\begin{multicols}{2}
\textbf{Connections:}
\begin{itemize}[itemsep=3pt]
  \item Lesson 2.4: DOFS described direction and strength in words --- $r$ assigns a number.
  \item Lesson 2.6: slope of regression line $= r \cdot \dfrac{s_y}{s_x}$.
  \item Lessons 2.8--2.9: $r^2$ measures how much variability in $y$ is explained by $x$.
\end{itemize}

\columnbreak

\textbf{Reflection:}\\
Why must you \emph{always} look at the scatterplot even after computing $r$?\\[2pt]
\writeline\\[1pt]\writeline\\[1pt]\writeline
\end{multicols}
\end{spiralbox}

\vspace{0.08in}

\begin{notesbox}{Additional Notes}
\vspace{0pt}
\writeline\\\writeline\\\writeline\\\writeline\\\writeline
\end{notesbox}

\end{document}
